% =====================================================================
% Section — Challenges: where state-of-the-art native implementations
% reach the limit of the instruction set, and what microcode offers
% beneath it. Motivates the two worked examples (Keccak, Curve25519).
%
% Intended placement: after Section 3 (the microcode layer) and before
% Section 4 (Keccak), so that each worked example reads as the removal
% of a named, ISA-level obstacle established here.
%
% FACTS verified for this draft (2026-06-28):
%   - Goldmont (Celeron N3350) is OUT-OF-ORDER, 3-wide, 3 integer ALU
%     ports (Wikipedia Goldmont; do NOT call it in-order).
%   - Goldmont's ISA reaches SSE4.2/AES/SHA and includes NEITHER BMI2
%     (MULX) NOR ADX (ADCX/ADOX) NOR AVX/AVX2. Verified two ways:
%     `gcc -march=goldmont -dM -E` enables only AES/POPCNT/SHA/SSE4.2;
%     Wikipedia Goldmont lists no BMI/ADX/AVX. ADX debuted on Broadwell.
%   - 5x51 schoolbook = 25 partial products; 4x64 saturated = 16.
% Citation keys: those already used in sec3-5 are reused; new keys are
% listed with ready-to-paste BibTeX at the end of the handoff message.
% =====================================================================

\section{Challenges: The Ceilings of the Instruction Set}
\label{sec:challenges}

The implementations we set out to surpass are not naive. For each of the two
primitives studied in this work the state of the art on commodity x86 is the
product of sustained expert effort: hand-written assembly by the primitive's own
designers~\cite{PKC:Bernstein06,bernsteinHighSpeedHighSecurity2012}, a formally
verified arithmetic generator \cite{erbsenSimpleHighLevel2019}, and a dedicated
cryptographic superoptimiser \cite{kuepperCryptOpt2023}. What unites them is that every one is expressed in
the x86-64 instruction set, and it is the instruction set itself, rather than any
shortfall of engineering, that bounds them. Two of its properties set those
bounds, and neither can be relaxed from above the architectural interface at
which a compiler or an assembly programmer works: the size of the
general-purpose register file, and the presence of exactly one architectural
carry flag. The first binds Keccak; the second binds the Curve25519 field
arithmetic. Both are ceilings beneath which the microcode layer of
\Cref{sec:ucode-layer} operates, and the purpose of this section is to state each
precisely, so that the two results that follow may be read as the removal of a
named obstacle rather than as an unexplained claim of speed.

The target throughout is a single Goldmont core, the \hbox{Celeron~N3350}, and
two of its characteristics matter to what follows. It is an out-of-order,
three-wide superscalar core with three integer execution pipelines, so the usual
intuition that a register-poor or dependency-bound program will simply be
re-ordered into the idle units is available to it and, as we shall see, still
does not rescue either primitive. And its instruction set extends only as far as
SSE4.2, AES and SHA: it possesses neither the AVX2 vector width nor the BMI2
\texttt{MULX} and ADX \texttt{ADCX}/\texttt{ADOX} instructions that Intel
introduced on later, larger cores \cite{intelLargeIntegerArithmetic2012}, ADX
having first appeared with Broadwell. The two escape hatches that the state of
the art relies on elsewhere --- vectorisation for Keccak and parallel
carry-chains for the field arithmetic --- are therefore both closed on this part,
which is what makes its limits a clean test of the layer beneath it.

% ---------------------------------------------------------------------
\subsection{Keccak: A State That Does Not Fit}
\label{sec:challenge-keccak}

\paragraph{The register-file ceiling.}
A Keccak-$f[1600]$ state is 1600 bits, laid out as 25 lanes of 64 bits, and the
permutation transforms it through 24 rounds of $\theta,\rho,\pi,\chi,\iota$
\cite{bertoniKeccakReference2011,technologySHA32015}. The x86-64 architecture exposes
sixteen general-purpose registers, one of which is the stack pointer, leaving on
the order of fifteen to hold data. The state thus needs twenty-five registers
where the architecture offers fifteen, so that a single round --- which reads
every lane once in the column parity of $\theta$ and again in $\chi$ --- cannot
keep even two-thirds of its lanes resident. Native code is obliged to spill the
majority of the state to the first-level cache and to reload it, and it does so on
each of the 24 rounds. The quantity that dominates a scalar Keccak permutation on
a core of this class is in consequence not arithmetic but the movement of state
between the registers and L1.

\paragraph{What the state of the art does, and why it does not lift the ceiling.}
The fastest scalar implementations attack the shortage in two ways, both legible
in the very names of the SUPERCOP contenders we measure against in
\Cref{sec:keccak-perf}. They unroll all 24 rounds, so that the register allocator
can carry a lane's value across a round boundary instead of reloading it, and they
apply lane complementing to eliminate some of the bitwise complements of $\chi$
that would otherwise cost an instruction and a register apiece, both of which are
distributed in the eXtended Keccak Code Package \cite{XKCP}; the contender
\texttt{opt64lcu24} announces them, in its lane-complement ``\texttt{lc}'' and its
unrolled-24 ``\texttt{u24}''. Neither device removes the ceiling. The unrolled
body is still register-starved within each round; full unrolling exchanges the
spill for a large static instruction footprint; and the in-place update of
$\chi$, in which writing one output lane destroys an input that a later lane of
the same row still requires, forces auxiliary temporaries and the moves that fill
them precisely when registers are scarcest. The remaining escape, vectorisation,
does not apply to the case on the critical path: a SIMD Keccak parallelises
several \emph{independent} permutations across the lanes of a vector register,
which accelerates a batched or tree-structured hash but does nothing for the
single permutation that a sequential hash must compute, and in any event
Goldmont's vector unit reaches only SSE4.2 and lacks the AVX2 width those
implementations are written for. On this microarchitecture a single Keccak
permutation has, within the instruction set, no route to register residency.

\paragraph{The opportunity.}
The obstacle is one of the \emph{architectural} register file, not of the
algorithm: the permutation refers to only 25 lanes and a handful of temporaries,
comfortably within the 32-entry physical register file that the microcode layer
exposes (\Cref{sec:ucode-layer}). The whole state can therefore be loaded once
and held resident across all 24 rounds, so that the register-to-memory traffic is
paid a single time per permutation rather than once per round, and the recovered
registers can be spent on the per-round transformations --- above all a $\chi$
that needs no auxiliary moves --- that register-starved native code cannot afford.
This is the ceiling that \Cref{sec:keccak} removes, and the reason its advantage
is structural rather than a matter of faster individual operations.

% ---------------------------------------------------------------------
\subsection{Curve25519: A Carry Chain Through a Single Flag}
\label{sec:challenge-curve25519}

\paragraph{Field arithmetic and the cost of its carries.}
A X25519 scalar multiplication is dominated by multiplications and squarings in
the field $\mathrm{GF}(2^{255}-19)$, of which a Montgomery ladder performs on the
order of a thousand \cite{PKC:Bernstein06}. Each is a multi-precision operation
whose cost on a 64-bit core is set as much by the propagation of carries between
machine words as by the word multiplications themselves. The x86-64 architecture
exposes a single carry flag, and the add-with-carry instruction from which a
multi-precision accumulation is built both reads and writes that one flag. A
sequence of such additions is therefore a strict dependency chain, each addition
waiting on the carry the previous one produced; and because that dependency is
genuine, threaded through the one shared flag, an out-of-order core cannot overlap
it however many integer units stand idle. The three execution pipelines of
Goldmont do not help here, for the limit is a true data dependency rather than a
scarcity of units.

\paragraph{The instructions the chip does not have.}
Intel recognised this bottleneck and addressed it in hardware. The ADX extension
adds \texttt{ADCX} and \texttt{ADOX}, which maintain two independent carry chains
on the carry and overflow flags respectively, while BMI2 adds \texttt{MULX},
which forms a $64\times64$ product without disturbing the flags; together they let
a multi-precision multiply advance two carry chains at once and overlap them with
the products that feed them \cite{intelLargeIntegerArithmetic2012}. The fastest
saturated-radix field implementations on current Intel parts are built directly on
these instructions, among them the \texttt{maa4}, \texttt{maa5} and \texttt{maax}
field-arithmetic variants distributed in lib25519
\cite{nathSarkarEfficientArithmetic2022,kuepperCryptOpt2023}.
They are, however, unavailable on our target: Intel introduced ADX only with
Broadwell, and Goldmont's instruction set, reaching SSE4.2, AES and SHA, includes
neither \texttt{MULX} nor \texttt{ADCX}/\texttt{ADOX}. On this microarchitecture
native code has one carry flag and hence one carry chain, and the standard
hardware remedy for serialised carries is simply absent.

\paragraph{The representation dilemma.}
The way native implementations cope is to change the number representation, and
the choice is a true dilemma rather than a free win. A saturated representation
packs the 255-bit element into four 64-bit limbs, the choice of the
\texttt{amd64-64} implementation \cite{bernsteinHighSpeedHighSecurity2012}, which
minimises the number of word multiplications --- sixteen partial products for the
$4\times4$ schoolbook --- but generates full-width carries between limbs, exactly
the long carry chains that want ADX. A reduced-radix representation instead
spreads the element across five 51-bit limbs, the choice of \texttt{amd64-51}
\cite{bernsteinHighSpeedHighSecurity2012} and of fiat-crypto, leaving thirteen
spare bits in each 64-bit word; carries are then short and need not be resolved
immediately, which is what makes the form attractive on a core without ADX, but
the extra limb raises the multiplication count to twenty-five
\cite{PKC:Bernstein06,erbsenSimpleHighLevel2019}. Neither representation escapes
the instruction set: the saturated form is fast only with carry instructions
Goldmont lacks, and the reduced form pays for its lazy carries in additional
multiplications. The verified generator, the superoptimiser, and the designers'
own assembly that we compare against
\cite{erbsenSimpleHighLevel2019,kuepperCryptOpt2023,bernsteinHighSpeedHighSecurity2012}
all operate within these two options and within the single-flag limit they share.
Nor is there a vectorised escape: the fastest vectorised X25519, Sandy2x
\cite{chouSandy2x2015}, is built on the AVX instruction set and therefore does not
run on Goldmont at all, exactly as the SIMD Keccak of
\Cref{sec:challenge-keccak} does nothing for the single permutation on the
critical path.

\paragraph{The opportunity, and its boundary.}
The microcode layer is not bound by the one architectural flag. The findings of
\Cref{sec:findings} establish a per-register condition that each accumulator
carries independently, so that several carry chains may advance at once on private
flags --- the very capability ADX would have supplied in hardware, here expressed
beneath the instruction set on a part that does not provide it. In the
reduced-radix form this lets a field multiplication overlap its carry propagation
with its multiplications in a manner no native code on this core can match, which
is the obstacle that \Cref{sec:curve25519} removes. The same finding also fixes
the boundary of the advantage: the architectural flag is frozen on entry to a
patch and cannot be threaded through the long, full-width carry chain that a
saturated $4\times64$ multiplication produces, so the representation that suits the
native core best is the one that suits the layer worst --- a tension that
\Cref{sec:curve25519} confronts directly when it reports where the microcode
ceases to win.

% ---------------------------------------------------------------------
\medskip
\noindent Both ceilings are properties of the architectural interface --- the
extent of the register file and the width of the carry machinery --- and neither
yields to a better compiler or a more patient assembly programmer, because both
are fixed below the level at which that effort is spent. Microcode is fixed lower
still. The two sections that follow take each ceiling in turn: \Cref{sec:keccak}
shows residency defeating the register-file limit on Keccak, and
\Cref{sec:curve25519} shows concurrent private-flag carry chains defeating the
single-flag limit on the Curve25519 field arithmetic, each stated together with
the point at which its advantage ends.
